\documentclass[11pt]{article}
\usepackage[margin=1in]{geometry}
\usepackage{amsmath, amssymb}
\usepackage{booktabs}
\usepackage{hyperref}
\usepackage{xcolor}
\usepackage{graphicx}

\title{Expected Improvements from PRO Depth Data Cleaning\\
       \large Bark-Brown-02 Box-Family Stereo Ablation}
\author{Jannick Sanchez}
\date{\today}

\begin{document}
\maketitle

\section{Background}

The stereo depth refinement models (MVStereoUNet, MVStereoDINOUNet) take
PRO depth maps as their primary input signal and learn to refine them
toward ground-truth (GT) Blender depth.  The quality of this refinement
depends directly on PRO depth being a meaningful, structured signal at
training time.

\section{The Corruption Problem}

PRO depth maps are stored as \texttt{.npy} files on disk.  An audit of
the 100 available tree directories found two categories of corrupt files:

\begin{table}[h!]
\centering
\begin{tabular}{lrr}
\toprule
Corruption type & Count & Cause \\
\midrule
Empty (0-byte)    & 57 & failed write / truncated transfer \\
Stub (128-byte)   & 43 & NumPy header only, no array data \\
\midrule
Total crash-causing & 100 & raise \texttt{EOFError} on load \\
\bottomrule
\end{tabular}
\caption{Corrupt PRO depth files by type.}
\end{table}

These affect \textbf{69 of 100 trees} (69\%) — at least one corrupt file
per tree.  A catch block in the dataset loader (\texttt{except (ValueError,
OSError)}) did not cover \texttt{EOFError}, so these files crashed the data
loader.

\subsection*{Fix Applied}

The exception tuple was extended to include \texttt{EOFError}:
\begin{verbatim}
except (ValueError, OSError, EOFError) as e:
    # fall back to constant depth = _PRO_BETA
\end{verbatim}

When triggered, the loader substitutes a constant depth tensor
(\texttt{\_PRO\_BETA} $\approx 1.556$ m, the median-calibration offset) for
the missing PRO map.

\section{Effect of Constant-Fallback on Training}

\subsection*{What the model receives}

For any sample drawn from a corrupt tree, the model's input is:
\[
  d_{\text{pro}}^{(i)} = \mathbf{1}_{H \times W} \cdot \beta_{\text{PRO}}
  \quad \text{(constant spatial map)}
\]
instead of the actual depth structure.  This means:

\begin{enumerate}
  \item \textbf{No depth gradient in input.}  Edges, foreground/background
    transitions, and structural cues the model is supposed to refine are
    entirely absent.
  \item \textbf{Trivial regression target.}  The model must predict
    varying GT depth from a flat input — this degrades to a pure
    hallucination task, which is unrelated to the intended refinement task.
  \item \textbf{Conflicting supervision.}  Within a mini-batch, the model
    receives both valid PRO depth samples (structured input $\to$ precise GT)
    and fallback samples (flat input $\to$ varying GT).  The gradient signal
    from fallback samples actively contradicts the signal from clean samples.
\end{enumerate}

\subsection*{Scale of contamination under random sampling}

With 69/100 trees having at least one bad PRO file, a random 8-tree
training set has a probability of:
\[
  P(\text{all 8 clean}) = \frac{\binom{31}{8}}{\binom{100}{8}} \approx 0.5\%
\]
of being fully clean.  In practice, no random seed in $[1, 99{,}999]$
produces a completely clean 8-tree training split.

\section{Clean Seed 6 — Data Selection Strategy}

Rather than relying on random sampling, seed 6 builds manifests by
\emph{restricting the sampling pool} to clean trees only:

\begin{enumerate}
  \item Parse \texttt{bad\_pro\_files.txt} to identify the 69 affected trees.
  \item Filter the box-family pool to the remaining \textbf{24 clean train
    trees} and \textbf{7 clean val trees}.
  \item Sample 8 train + 2 val trees from these clean pools using
    \texttt{random.seed(6)}.
\end{enumerate}

Selected trees (seed 6):
\begin{itemize}
  \item \textbf{Train:} \texttt{lpy\_envy\_00000, 00002, 00003, 00005,
    00024, 00053, 00066, 00079}
  \item \textbf{Val:} \texttt{lpy\_envy\_00070, 00091}
\end{itemize}

Every sample in the resulting manifests (240 train rows, 60 val rows) has a
valid, structured PRO depth map — the constant fallback is never triggered.

\section{Expected Improvements}

\subsection*{Training signal quality}

With clean data, every gradient step is computed from a meaningful
(input structure $\to$ GT refinement) pair.  The per-sample loss
gradient directly reflects depth refinement progress rather than
hallucination from flat inputs.

\subsection*{Convergence}

Gradient noise from conflicting fallback samples is eliminated.
We expect:
\begin{itemize}
  \item Lower training loss at early epochs (fewer misleading samples
    to overcome).
  \item More stable validation curve — validation loss should decrease
    monotonically rather than oscillating when clean/corrupt batches
    alternate.
  \item Earlier plateau / earlier stopping under patience=10.
\end{itemize}

\subsection*{Ablation validity}

The cross-view fusion ablation (fusion vs.\ no-fusion) compares:
\begin{align*}
  \text{With fusion:}\quad    &\hat{d} = \text{Decode}(\text{FuseConv}([b_0, b_1]),\; s_i) \\
  \text{Without fusion:}\quad &\hat{d} = \text{Decode}(b_i,\; s_i)
\end{align*}

Under contaminated data, fallback samples provide identical bottlenecks
for both views ($b_0 = b_1$ since both inputs are the same constant map),
artificially inflating the apparent benefit of fusion.  With clean data,
the two views carry genuinely different depth structure, making the
fusion comparison meaningful.

\subsection*{GT depth quality}

The GT Blender depth maps are rendered synthetically and are not
themselves corrupt.  However, samples whose PRO depth files were missing
may also have had other data inconsistencies (e.g., missing mask files
recorded in the same corrupt write event).  Excluding these trees reduces
the risk of subtle label noise in the training set.

\section{Summary Table}

\begin{table}[h!]
\centering
\begin{tabular}{lcc}
\toprule
Property & Seeds 1--5 (random) & Seed 6 (clean) \\
\midrule
Fraction of trees with any bad PRO file & 69\% & 0\% \\
Constant-fallback samples in training   & $\sim$30--50\% & 0\% \\
Gradient signal per step                & Mixed & Clean \\
Fusion ablation confounded by flat inputs & Yes & No \\
Training trees & 8 (random pool) & 8 (clean pool of 24) \\
Val trees & 2 (random pool) & 2 (clean pool of 7) \\
\bottomrule
\end{tabular}
\caption{Comparison of random (seeds 1--5) vs.\ clean (seed 6) data splits.}
\end{table}

\section{Limitations}

\begin{itemize}
  \item Only 24 clean train trees are available for the box-family.  The
    effective training pool is smaller than the full 80-tree pool, which
    may reduce tree-level diversity.
  \item A single clean seed provides one data point, not a distribution.
    Comparing seed 6 to the mean of seeds 1--5 may partly reflect tree
    selection bias rather than data cleaning benefit.
  \item The clean seed result should be interpreted as an \emph{upper bound}
    on what the corrupted runs could have achieved, not as a definitive
    performance estimate.
\end{itemize}

\end{document}
